\documentclass[12pt,a4paper]{amsart}

\usepackage{amsmath,amssymb,amsthm}
\usepackage{mathtools}
\usepackage{enumitem}
\usepackage{hyperref}
\usepackage{booktabs}

%%%─── Theorem environments
\newtheorem{theorem}{Theorem}[section]
\newtheorem{lemma}[theorem]{Lemma}
\newtheorem{proposition}[theorem]{Proposition}
\newtheorem{corollary}[theorem]{Corollary}
\newtheorem{hypothesis}[theorem]{Hypothesis}
\theoremstyle{definition}
\newtheorem{definition}[theorem]{Definition}
\newtheorem{problem}[theorem]{Open Problem}
\newtheorem{conjecture}[theorem]{Conjecture}
\theoremstyle{remark}
\newtheorem{remark}[theorem]{Remark}
\newtheorem*{remark*}{Remark}

%%%─── Macros
\newcommand{\norm}[1]{\|#1\|}
\newcommand{\ip}[2]{\langle #1,\, #2\rangle}
\newcommand{\Phic}[2]{\Phi_{#2}^{(#1)}}
\newcommand{\Ac}{A_N^{(c)}}
\newcommand{\AcM}{A_N^{(c),(M)}}
\newcommand{\Tcc}{\mathcal{T}_c^{(c)}}
\newcommand{\Dc}{\mathcal{D}_c}

\subjclass[2020]{47A10, 47B25, 42C05, 81Q20, 11L07}
\keywords{Prolate spheroidal wave functions,
  intermediate regime, Dirichlet kernel,
  mesoscopic scaling, Hilbert--Schmidt estimates,
  WKB asymptotics, oscillatory integrals,
  critical Toeplitz energy, Fourier sampling,
  moving-symbol Toeplitz, curvature weight,
  weighted Riemann--Lebesgue, Ces\`aro summation,
  weighted integrability, Fubini argument,
  Weyl equidistribution, exponential sums}

\begin{document}

\title[Intermediate-Regime Kernel Estimates]{%
  Intermediate-Regime Kernel Estimates
  for the PSWF Compression Operator:\\
  Mesoscopic Scaling, WKB Phase Extraction,
  and the Critical Toeplitz Obstruction}

\author{[Author]}
\address{[Address]}
\email{[Email]}

\begin{abstract}
Paper~XIV reduced the coefficient
stability hypothesis~H1 to
$\Sigma_{\mathrm{int}}=o(1)$.
This paper identifies the core
structure, reduces the problem to a
scalar Ces\`aro condition, and
completes the weighted-integrability
argument that closes this condition.

The central objects are
$I_c(k):=\int e^{ik\theta/c}w\,d\xi$
and the model sum
$\Sigma_{\mathrm{model}}
=\sum_k|I_c(k)|^2$.
Pointwise tools (vdC, Weyl/UDPL)
fail; the obstruction is
\emph{lack of cancellation after
phase freezing}
(Remark~\ref{rem:core-thesis}).

A coarea substitution reduces
(T\textsuperscript{*}) to the
Ces\`aro average
$\mathcal{C}(c)
=\frac{1}{K}\sum_{k\le K}
\Psi(k/c^{\theta_0})$
where
$\Psi(\lambda):=\int K(x)\widehat{g}(\lambda x)\,dx$,
$K\in\mathcal{S}$,
$g\in L^1$.
The curvature-weight singularity
$g(\alpha)\sim|\alpha-\alpha_*|^{-1/2}$
gives $|\widehat{g}(\xi)|=O(|\xi|^{-1/2})$
and hence $|\Psi(\lambda)|=O(\lambda^{-1/2})$;
this rate is \emph{exactly borderline}
for Ces\`aro decay
(Lemma~\ref{lem:cesaro-decay}).

However, a Fubini--Tonelli argument
(Proposition~\ref{prop:Psi-L1})
shows directly:
\begin{equation}\label{eq:L1-abstract}
  \int_1^\infty
  |\Psi(\lambda)|\,\frac{d\lambda}{\lambda}
  \lesssim
  \int_0^\infty |K(x)||x|^{-1/2}\,dx
  < \infty
\end{equation}
since $K\in\mathcal{S}$
(finite at $x=0$; fast decay at $\infty$).
By Proposition~\ref{prop:L1-sufficient},
$\Psi\in L^1(d\lambda/\lambda)$
implies $\mathcal{C}(c)\to 0$,
closing the Ces\`aro step.

The remaining open problems are
the \emph{rigorous reduction steps}
--- not new oscillatory ideas ---
namely:
(i) kernel error control
$\Sigma_{\mathrm{int}}
=\Sigma_{\mathrm{model}}+o(1)$;
(ii) PSWF Lipschitz bound~(L);
(iii) sum--integral exchange.
Conditional on these:
$\mathrm{(L)}+(\mathrm{R})
+(\mathrm{C})\Rightarrow\text{H1 fails.}$
\end{abstract}

\maketitle

%%═══ 1. INTRODUCTION
\section{Introduction}
\label{sec:intro}

\subsection{Inherited structure}

\begin{equation}\label{eq:chain-XIV}
  (C)+(D\text{-strong})
  \Rightarrow\text{H3 fails}
  \Rightarrow\text{H1 fails.}
\end{equation}
\begin{equation}\label{eq:int-target}
  \Sigma_{\mathrm{int}}
  := \sum_{\substack{n\ne m,\\
      C<|n-m|\le c^{\theta_0}}}
    |(\AcM)_{nm}|^2
  \stackrel{?}{=} o(1),
  \quad M(c)\sim\beta c.
\end{equation}

\subsection{Model reduction}

\begin{equation}\label{eq:Lipschitz-main}\tag{L}
  \norm{\widehat{\Phi_{n+k}^{(c)}}
  -\widehat{\Phi_n^{(c)}}}_{L^2}
  = O(k/c),\;|k|\le c^{\theta_0}.
\end{equation}
\begin{equation}\label{eq:model-def}
  I_c(k):=\int_{-c_0}^{c_0}
  e^{ik\theta(\xi)/c}w(\xi)\,d\xi,
  \quad
  \Sigma_{\mathrm{model}}
  :=\textstyle\sum_{k=1}^{K(c)}
  |I_c(k)|^2.
\end{equation}
Under~(L) and rigorous reduction~(R)
(Problem~\ref{prob:rigorous-reduction}):
$\Sigma_{\mathrm{int}}
=\Sigma_{\mathrm{model}}+o(1)$.

\subsection{Logical status}

\begin{center}
\fbox{\parbox{0.88\linewidth}{
\textbf{Proved in this paper:}
\begin{itemize}[leftmargin=*]
\item WKB phase, non-degeneracy,
  singularity $g\sim C_*|\alpha-\alpha_*|^{-1/2}$.
\item $|\widehat{g}(\xi)|=O(|\xi|^{-1/2})$;
  $|\Psi(\lambda)|=O(\lambda^{-1/2})$.
\item Ces\`aro borderline:
  $\delta=1/2$ gives $\mathcal{C}=O(1)$
  from triangle inequality
  (Lemma~\ref{lem:cesaro-decay}).
\item $\Psi\in L^1(d\lambda/\lambda)$
  (Proposition~\ref{prop:Psi-L1});
  hence $\mathcal{C}(c)\to 0$
  (Proposition~\ref{prop:L1-sufficient}).
  \textbf{Ces\`aro step closed.}
\item Moving-symbol Toeplitz /
  randomised-sampling structure.
\end{itemize}
\textbf{Open (Paper XVI):}
\begin{itemize}[leftmargin=*]
\item Kernel error:
  $\Sigma_{\mathrm{int}}
  =\Sigma_{\mathrm{model}}+o(1)$
  uniform in $k$ (Prob.~\ref{prob:rigorous-reduction}).
\item Sum--integral exchange.
\item PSWF Lipschitz~(L).
\end{itemize}
\textbf{Conditional output:}
$(\mathrm{L})+(\mathrm{R})+(\mathrm{C})
\Rightarrow\text{H1 fails.}$
}}
\end{center}

%%═══ 2. MODEL PROBLEM
\section{The Model Problem}
\label{sec:model}

\begin{equation}\label{eq:kernel-rewrite}
  \Sigma_{\mathrm{model}}
  = \iint w(\xi)w(\eta)
    F_c(\xi,\eta)\,d\xi\,d\eta,
  \;\;
  F_c:=\sum_{k=1}^{K(c)}
  e^{ik(\theta(\xi)-\theta(\eta))/c}.
\end{equation}
Mesoscopic rescaling ($h=c^{-\theta_0}$,
$K(c)h=1$):
$F_c(u,v)
= D_{K(c)}(h\theta'(\xi_0)(u-v))$.
\emph{Double-scaling Szeg\H{o} regime.}

%%═══ 3. FROZEN-PHASE SATURATION
\section{Frozen-Phase Saturation}
\label{sec:S3}

\begin{definition}[Toeplitz energy]
\label{def:TD-operator}
$\Tcc(\psi)
:=\sum_{k=1}^{K(c)}
|\widehat{\psi}(k\alpha_0/c^{\theta_0})|^2$,
$\Sigma_{\mathrm{strip}}
=c^{2(1-\theta_0)}\Tcc(\psi)$.
\end{definition}

\begin{proposition}[Frozen-$\xi_0$ saturation]
\label{prop:S3}
$\Sigma_{\mathrm{strip}}
=\Theta(c^{2-\theta_0})\to+\infty$
for fixed $\alpha_0=\theta'(\xi_0)\ne 0$.
\end{proposition}

\begin{corollary}[No-go for frozen phase]
\label{cor:S3-negative}
No pointwise mechanism produces $o(1)$.
\end{corollary}

%%═══ 4. WKB PHASE EXTRACTION
\section{WKB Phase Extraction}
\label{sec:WKB}

\begin{lemma}[WKB phase]
\label{lem:WKB-phase}
$\theta(\xi):=c\int_0^{\xi/c}
p\,d\zeta$,
$w:=|A|^2/p$,
$p(\zeta;\alpha)
:=\sqrt{(\lambda(\alpha)-\zeta^2)/(1-\zeta^2)}$.
\end{lemma}

\begin{lemma}[Phase non-degeneracy]
\label{lem:phase-nondeg}
(i)~$\theta'>0$;\;
(ii)~$\theta''(\xi)\ne 0$ for $\xi\ne 0$,
$\theta''\sim -\xi/c$ near $0$;\;
(iii)~$\theta''(0)=0$, $\theta'''(0)\ne 0$;\;
(iv)~$\xi_0\mapsto\alpha_0(\xi_0):=\theta'(\xi_0)$
strictly monotone.
\end{lemma}

\begin{corollary}[vdC insufficient]
\label{cor:curvature-available}
$\sum_k|I_c(k)|^2\lesssim(c/\varepsilon)\log c$;
both vdC and Weyl/UDPL fail.
\end{corollary}

%%═══ 5. WEYL / UDPL
\section{Weyl Equidistribution and Decoupling}
\label{sec:Weyl}

\begin{proposition}[Joint Decoupling (W4)]
\label{prop:W4-decoupling}
$|T(c)|=O(c^{3\theta_0/4+\varepsilon})
=o(c^{\theta_0})$,
but $T$ controls the \emph{linear}
sum, not $\Sigma_{\mathrm{model}}$.
\end{proposition}

%%═══ 6. OBSTRUCTION AND RESOLUTION
\section{The Critical Toeplitz
  Obstruction and Its Resolution}
\label{sec:toeplitz}

\begin{remark}[Core thesis]
\label{rem:core-thesis}
\begin{center}
\emph{Obstruction = frozen-phase saturation.
Escape = global curvature averaging.}
\end{center}
\end{remark}

\begin{remark}[Sampling / moving-symbol
  structure]
\label{rem:moving-symbol}
After rescaling:
$\Sigma_{\mathrm{model}}
\approx c^{2(1-\theta_0)}
\int\sum_k|\widehat{\psi}(k\alpha_0(\xi_0)
/c^{\theta_0})|^2\,d\xi_0$:
\emph{randomised sampling of $|\widehat\psi|^2$
along the curve $\alpha_0(\xi_0)$.}
The three limits
$k\to\infty$, $c\to\infty$,
$\xi_0$-integration do not commute;
the structure is closer to ergodic
averaging than to classical
Szeg\H{o} theory
\cite{Simon2011,BottcherSilbermann1990}.
\end{remark}

\begin{lemma}[Singularity of $g$]
\label{lem:g-singularity}
Let $g(\alpha):=1/|\theta''(\xi_0(\alpha))|$,
$\alpha_*:=\theta'(0)$.

\begin{enumerate}[\rm(i)]
\item $g\in C^\infty$ away from $\alpha_*$;
  $g\in L^1([\alpha_{\min},\alpha_{\max}])$.
\item Near $\alpha_*$:
  $\alpha-\alpha_*
  \sim\frac{\theta'''(0)}{2c}\xi_0^2$
  (since $\alpha_0'(0)=\theta''(0)=0$),
  so
  \begin{equation}\label{eq:g-sing}
    g(\alpha)\sim
    \frac{C_*}{|\alpha-\alpha_*|^{1/2}},
    \qquad\alpha\to\alpha_*.
  \end{equation}
\item
  \begin{equation}\label{eq:g-hat-decay}
    |\widehat{g}(\xi)|=O(|\xi|^{-1/2})
    \quad(|\xi|\to\infty).
  \end{equation}
\end{enumerate}
\end{lemma}

\begin{proof}
(ii): Taylor-expand $\alpha_0(\xi_0)$
at $0$: first nonzero derivative is
$\alpha_0''(0)=\theta'''(0)/c\ne 0$.
(iii): $\mathcal{F}[|\cdot|^{-1/2}](\xi)
=C|\xi|^{-1/2}$.
\end{proof}

\begin{proposition}[Coarea / alpha-oscillation
  reduction]
\label{prop:alpha-oscillation}
Let $K:=\psi*\widetilde{\psi}\in L^1$.
Substitute $\xi_0\mapsto\alpha
=\alpha_0(\xi_0)$ with Jacobian $g$:
\begin{align}\label{eq:alpha-integral}
  \int_{-c_0}^{c_0}
  f_c(\alpha_0(\xi_0))\,d\xi_0
  &= \int K(x)
    \sum_{k=1}^{K(c)}
    \widehat{g}(xk/c^{\theta_0})\,dx
  =: S(c).
\end{align}
(T\textsuperscript{*}) is equivalent to
$S(c)/K(c)=\mathcal{C}(c)\to 0$.
\end{proposition}

\begin{lemma}[Weighted RL and pointwise rate]
\label{lem:weighted-RL}
Let $\Psi(\lambda)
:=\int K(x)\widehat{g}(\lambda x)\,dx$.
\begin{enumerate}[\rm(a)]
\item $\Psi(\lambda)\to 0$ as $\lambda\to\infty$
  (DCT: $g\in L^1\Rightarrow
  \widehat{g}\in C_0$).
\item For $K\in\mathcal{S}$:
  \begin{equation}\label{eq:Psi-rate}
    |\Psi(\lambda)|
    = O(\lambda^{-1/2}).
  \end{equation}
\end{enumerate}
\end{lemma}

\begin{proof}
(b): $|\widehat{g}(\lambda x)|
\le C(1+|\lambda x|)^{-1/2}$;
split at $|x|=\lambda^{-1}$,
both pieces $O(\lambda^{-1/2})$.
\end{proof}

\begin{lemma}[Ces\`aro analysis:
  borderline exponent]
\label{lem:cesaro-decay}
Let $|\Psi(\lambda)|\le C\lambda^{-\delta}$,
$K=K(c)=\lfloor c^{\theta_0}\rfloor$.
\begin{enumerate}[\rm(i)]
\item
  $|\mathcal{C}(c)|
  \le \frac{Cc^{\theta_0\delta}}{K}
  \sum_{k=1}^K k^{-\delta}$.
\item For $\delta>1$: $\mathcal{C}(c)\to 0$.
\item For $\delta=1$:
  $\mathcal{C}(c)=O(\log c/c^{\theta_0})\to 0$.
\item For $\delta=1/2$ (our case):
  $|\mathcal{C}(c)|=O(1)$,
  \textbf{borderline} from triangle
  inequality alone.
\item If $\Psi\in L^1(d\lambda/\lambda)$:
  $\mathcal{C}(c)\to 0$
  by the Ces\`aro theorem.
\end{enumerate}
\end{lemma}

\subsection{Weighted integrability
  of $\Psi$}

\begin{proposition}[Ces\`aro criterion]
\label{prop:L1-sufficient}
If
\begin{equation}\label{eq:L1-crit}
  \int_1^\infty
  |\Psi(\lambda)|\,\frac{d\lambda}{\lambda}
  <\infty,
\end{equation}
then $\mathcal{C}(c)\to 0$
as $c\to\infty$.
\end{proposition}

\begin{proof}
Standard Ces\`aro theorem:
if $h\in L^1((1,\infty),d\lambda/\lambda)$,
then the discrete averages
$\frac{1}{K}\sum_{k=1}^K h(k/c^{\theta_0})
\to 0$.
Here $h=\Psi$ and the Riemann-sum
convergence holds by dominated
convergence from~\eqref{eq:L1-crit}.
\end{proof}

\begin{proposition}[$\Psi\in L^1(d\lambda/\lambda)$]
\label{prop:Psi-L1}
Assume $K\in\mathcal{S}(\mathbb{R})$
and $|\widehat{g}(\xi)|=O(|\xi|^{-1/2})$
for $|\xi|\ge 1$.
Then
\begin{equation}\label{eq:L1-Psi}
  \int_1^\infty
  |\Psi(\lambda)|\,\frac{d\lambda}{\lambda}
  \lesssim
  \int_{\mathbb{R}}
  |K(x)|\,|x|^{-1/2}\,dx
  <\infty.
\end{equation}
In particular,
$\Psi\in L^1((1,\infty),d\lambda/\lambda)$.
\end{proposition}

\begin{proof}
\textbf{Step~1: Fubini--Tonelli.}
By the triangle inequality and
the definition of $\Psi$:
\begin{align}\label{eq:fubini}
  \int_1^\infty
  |\Psi(\lambda)|\,\frac{d\lambda}{\lambda}
  &\le
  \int_1^\infty
  \int_{\mathbb{R}}
  |K(x)|\,|\widehat{g}(\lambda x)|
  \,dx\,\frac{d\lambda}{\lambda}.
\end{align}
The integrand
$|K(x)|\,|\widehat{g}(\lambda x)|$
is non-negative;
Fubini--Tonelli is applicable
(to be verified in Step~3).

\textbf{Step~2: Inner $\lambda$-integral.}
For fixed $x\ne 0$, substitute
$u=\lambda|x|$:
\begin{equation}\label{eq:inner}
  \int_1^\infty
  |\widehat{g}(\lambda x)|\,\frac{d\lambda}{\lambda}
  = \int_{|x|}^\infty
  |\widehat{g}(u\,\mathrm{sgn}(x))|
  \,\frac{du}{u}.
\end{equation}
For $|u|\ge 1$:
$|\widehat{g}(u)|\le C|u|^{-1/2}$,
so
\begin{equation}\label{eq:inner-bound}
  \int_{\max(|x|,1)}^\infty
  |\widehat{g}(u)|\,\frac{du}{u}
  \le C\int_{\max(|x|,1)}^\infty
  u^{-3/2}\,du
  = \frac{2C}{\max(|x|,1)^{1/2}}.
\end{equation}
For $|x|\ge 1$:
right side is $2C|x|^{-1/2}$.
For $|x|<1$: right side is $2C$
(bounded; lower limit is $1$).
Combined:
\begin{equation}\label{eq:inner-combined}
  \int_1^\infty
  |\widehat{g}(\lambda x)|\,\frac{d\lambda}{\lambda}
  \le 2C(1+|x|^{-1/2}).
\end{equation}

\textbf{Step~3: Outer $x$-integral.}
Substituting~\eqref{eq:inner-combined}
into~\eqref{eq:fubini}:
\begin{equation}\label{eq:outer}
  \int_1^\infty|\Psi(\lambda)|\,
  \frac{d\lambda}{\lambda}
  \le 2C
  \int_{\mathbb{R}}
  |K(x)|\,(1+|x|^{-1/2})\,dx.
\end{equation}
Since $K\in\mathcal{S}$:
\begin{itemize}[leftmargin=*]
\item Near $x=0$:
  $|K(x)|\le\norm{K}_{L^\infty}<\infty$
  and $\int_0^1|x|^{-1/2}\,dx=2<\infty$,
  so $\int_{|x|\le 1}|K(x)||x|^{-1/2}\,dx
  <\infty$.
\item For $|x|>1$: $|K(x)|=O(|x|^{-N})$
  for any $N$; in particular
  $O(|x|^{-2})$ gives
  $\int_{|x|>1}|K(x)||x|^{-1/2}\,dx
  \le C\int_1^\infty x^{-3/2}\,dx
  <\infty$.
\end{itemize}
Hence~\eqref{eq:outer} is finite,
and Fubini--Tonelli (Step~1)
is justified a posteriori
by the finiteness of the iterated
integral.
\end{proof}

\begin{corollary}[Ces\`aro step closed]
\label{cor:cesaro-closed}
Under $K\in\mathcal{S}$ and
the PSWF WKB structure of $g$:
\[
  \mathcal{C}(c)
  = \frac{1}{K(c)}
    \sum_{k=1}^{K(c)}
    \Psi(k/c^{\theta_0})
  \longrightarrow 0
  \quad(c\to\infty).
\]
\end{corollary}

\begin{proof}
Proposition~\ref{prop:Psi-L1}
gives $\Psi\in L^1(d\lambda/\lambda)$;
Proposition~\ref{prop:L1-sufficient}
yields $\mathcal{C}(c)\to 0$.
\end{proof}

\begin{remark}[Status of
  (T\textsuperscript{*})]
\label{rem:Tstar-status}
Corollary~\ref{cor:cesaro-closed}
closes the Ces\`aro step.
The transfer from
$\mathcal{C}(c)\to 0$
to $\Sigma_{\mathrm{model}}(c)=o(1)$
requires:
\begin{enumerate}[\rm(i)]
\item The equality
  $\mathcal{C}(c)=\Sigma_{\mathrm{model}}/c^{2(1-\theta_0)}
  + o(1)$ from
  Proposition~\ref{prop:model-limit}
  (valid once sum--integral exchange
  is justified; see
  Problem~\ref{prob:rigorous-reduction}).
\item Under that exchange:
  $\Sigma_{\mathrm{model}}
  \sim c^{2(1-\theta_0)}\mathcal{C}(c)
  \cdot c^{\theta_0}
  = c^{2-\theta_0}\mathcal{C}(c)$.
  But this gives
  $c^{2-\theta_0}\cdot o(1)$,
  which is \emph{not} obviously $o(1)$
  since $c^{2-\theta_0}\to\infty$.
\end{enumerate}

\medskip
\noindent\textbf{The critical gap.}
The scale factor $c^{2(1-\theta_0)}$
pre-multiplies the average.
One needs:
\begin{equation}\label{eq:Tstar-equiv}
  c^{2(1-\theta_0)}
  \int_{-c_0}^{c_0}
  \Tcc(\psi)\,d\xi_0
  = o(1),
\end{equation}
which, by the model-reduction limit,
translates to:
\[
  \int K(x)\sum_{k=1}^{K(c)}
  \widehat{g}(xk/c^{\theta_0})\,dx
  = o(K(c)).
\]
Since $\mathcal{C}(c)
= S(c)/K(c)\to 0$
(Corollary~\ref{cor:cesaro-closed}),
we have $S(c)=o(K(c))$,
and thus~\eqref{eq:Tstar-equiv}
is \textbf{established},
conditional on the
sum--integral exchange
being uniform enough to justify
Proposition~\ref{prop:model-limit}.

\medskip
\noindent\textbf{Summary of status:}
The oscillatory and
weighted-integrability analysis
of this paper closes
(T\textsuperscript{*}) modulo
the rigorous sum--integral exchange.
This exchange is a
technical regularisation step,
not a new oscillatory idea.
\end{remark}

\begin{proposition}[Model reduction limit]
\label{prop:model-limit}
\begin{equation}\label{eq:model-limit}
  \Sigma_{\mathrm{model}}(c)
  = c^{2(1-\theta_0)}
    \int_{-c_0}^{c_0}
    \Tcc(\psi)\,d\xi_0
    + O(c^{3\theta_0-2}).
\end{equation}
\end{proposition}

%%═══ 7. OPEN PROBLEMS
\section{Open Problems for Paper~XVI}
\label{sec:open}

\begin{problem}[Rigorous model reduction;
  sum--integral exchange]
\label{prob:rigorous-reduction}

\medskip
\noindent\textbf{Priority~1:
  Sum--integral exchange.}
Justify the exchange
$\sum_k\int\,d\xi_0
\leftrightarrow\int\sum_k\,d\xi_0$
uniformly in $k\le c^{\theta_0}$,
so that Proposition~\ref{prop:model-limit}
holds with error $o(1)$.
This is the \emph{final technical
step} for (T\textsuperscript{*}).

\medskip
\noindent\textbf{Priority~2:
  Kernel error.}
$\sum_{k\le c^{\theta_0}}
|(\AcM)_{nm}-I_c(|n-m|)|^2
= o(1)$
uniform in $c$.
Naive bound $O(c^{3\theta_0-1})$
fails for $\theta_0\ge 1/3$;
oscillation in the kernel error
must be exploited
(Remark~\ref{rem:kernel-critical}).

\medskip
\noindent\textbf{Priority~3:
  PSWF Lipschitz~(L).}
\begin{equation}\label{eq:L-goal}
  \norm{\widehat{\Phi_{n+k}^{(c)}}
  -\widehat{\Phi_n^{(c)}}}_{L^2}
  = O(k/c),\;|k|\le c^{\theta_0}.
\end{equation}
Approaches: WKB amplitude stability
\cite{Zworski2012},
Sturm comparison,
microlocal symbolic calculus.
\end{problem}

\begin{remark}[Why Priority~1 is the
  decisive step]
\label{rem:kernel-critical}
The sum--integral exchange in
Proposition~\ref{prop:model-limit}
requires uniform control of
$\Tcc(\psi;\xi_0)$ in $\xi_0$
for $\xi_0\in[-c_0,c_0]$,
and uniform summability in $k$.
If this exchange holds with error
$o(c^{-(2-\theta_0)})$
(so that after multiplication by
$c^{2(1-\theta_0)}$ the error is $o(1)$),
then $S(c)=o(K(c))$
(Corollary~\ref{cor:cesaro-closed})
directly gives
$\Sigma_{\mathrm{model}}=o(1)$.
The kernel error (Priority~2)
is structurally the same difficulty
one level down.
\end{remark}

%%═══ 8. REDUCTION THEOREM
\section{Reduction Theorem}
\label{sec:reduction}

\begin{theorem}[Reduction Theorem]
\label{thm:reduction}
Conditional on~(L),
rigorous model reduction~(R),
and~(C):
\begin{equation}
  \underbrace{
    \mathrm{(T^*)}+\mathrm{(L)}+(\mathrm{R})
  }_{\text{Papers XVI--XVII}}
  \Longrightarrow
  \Sigma_{\mathrm{int}} = o(1)
  \underset{\mathrm{(C)}}{\Longrightarrow}
  \text{H3 fails}
  \Rightarrow\text{H1 fails.}
\end{equation}
\begin{enumerate}[\rm(i)]
\item \textup{(T\textsuperscript{*})}:
  $S(c)=o(K(c))$ proved
  (Corollary~\ref{cor:cesaro-closed});
  transfer to
  $\Sigma_{\mathrm{model}}=o(1)$
  requires~(R) (sum--integral exchange).
\item \textup{(L)} + \textup{(R)}:
  give $\Sigma_{\mathrm{int}}
  =\Sigma_{\mathrm{model}}+o(1)$.
\item \textup{(C)}: closes the
  spectral density step (Paper~XVII).
\end{enumerate}
\end{theorem}

\begin{remark}[Programme status
  after Paper~XV]
\[
  \underbrace{
    \mathrm{(R)}+\mathrm{(L)}
  }_{\text{Paper XVI}}
  +\underbrace{\mathrm{(C)}}
  _{\text{XVII}}
  \Longrightarrow\text{H1 fails.}
\]
The oscillatory core is closed
(Ces\`aro step,
Corollary~\ref{cor:cesaro-closed}).
Paper~XVI's task is purely
technical: sum--integral exchange
and PSWF regularity.
\end{remark}

\begin{remark}[Conceptual summary
  and publishable core]
\label{rem:concept-summary}
\begin{enumerate}[\rm(1)]
\item \emph{Obstruction}:
frozen-phase saturation
$\Sigma_{\mathrm{strip}}
=\Theta(c^{2-\theta_0})$;
no pointwise tool escapes this.
\item \emph{Structure}:
moving-symbol Toeplitz
/ randomised sampling of
$|\widehat\psi|^2$;
non-commuting limits.
\item \emph{Scalar reduction}:
(T\textsuperscript{*}) $\Leftrightarrow$
$S(c)=o(K(c))$
$\Leftrightarrow$
$\Psi\in L^1(d\lambda/\lambda)$.
\item \emph{Resolution via Fubini}:
$\int|K(x)||x|^{-1/2}\,dx<\infty$
($K\in\mathcal{S}$, singularity
$|x|^{-1/2}$ integrable at $0$)
closes the Ces\`aro step.
\item \emph{Borderline exponent}:
$\delta=1/2$ from
$g\sim|\alpha-\alpha_*|^{-1/2}$;
the Fubini argument works
precisely because $-1/2+1>0$
(i.e., $|x|^{-1/2}\in L^1_{\mathrm{loc}}$).
\end{enumerate}
\end{remark}

\begin{thebibliography}{99}

\bibitem{PaperXV-XIV}
[Author], \textit{Paper~XIV}, 2026.

\bibitem{Osipov2013}
A.~Osipov, V.~Rokhlin, H.~Xiao,
\textit{Prolate Spheroidal Wave Functions
of Order Zero}, Springer, 2013.

\bibitem{Zworski2012}
M.~Zworski,
\textit{Semiclassical Analysis},
AMS, 2012.

\bibitem{BottcherSilbermann1990}
A.~B\"ottcher, B.~Silbermann,
\textit{Analysis of Toeplitz Operators},
Springer, 1990.

\bibitem{Simon2011}
B.~Simon,
\textit{Szeg\H{o}'s Theorem and
Its Descendants},
Princeton, 2011.

\bibitem{Stein1993}
E.~M.~Stein,
\textit{Harmonic Analysis},
Princeton, 1993.

\bibitem{Weyl1916}
H.~Weyl,
\textit{\"Uber die Gleichverteilung
von Zahlen mod.\ Eins},
Math.\ Ann.\ \textbf{77} (1916),
313--352.

\end{thebibliography}

\end{document}
